%% macros-care.tex : mises en évidence de la relecture assistée (CARE), chargé par
%% macros.tex s'il est présent. Ce fichier ne part jamais chez les tiers : make publier
%% écarte tout fichier de lib/ dont le nom contient « care », retire la ligne qui le
%% charge de la copie de macros.tex, et refuse de publier une source qui utilise ces
%% macros. Dans un document diffusé, il n'en reste donc aucune trace.

%% Boîte des mises en évidence CARE. Couleurs pensées pour la
%% protanopie : l'axe bleu-jaune reste discriminant (jaune = \IMPORTANT
%% d'Olivier, bleu = CARE), et la distinction ne repose pas que sur la
%% couleur (cadre + bandeau étiqueté, lisibles même en niveaux de gris).
\newtcolorbox{maboiteimportantcare}[2][]{colback=cyan!10!white, colframe=blue!60!black,fonttitle=\bfseries, colbacktitle=blue!70!black,enhanced,
attach boxed title to top center={yshift=-2mm},
  title={#2},#1}

%% Pendant de \IMPORTANT, réservé aux mises en évidence CARE ;
%% \IMPORTANT reste réservé aux annotations d'Olivier. Comme \IMPORTANT,
%% redevient du texte simple en mode DIFFUSIONFINALE.
\newcommand\IMPORTANTCARE[2][Important (CARE)]{
\ifnotmodeDIFFUSIONFINALE{
\begin{maboiteimportantcare}{#1}
#2
\end{maboiteimportantcare}
}
\ifmodeDIFFUSIONFINALE{#2}
}

%% \signeCARE[date]{Nom} : la signature d'un assistant au bas de ce qu'il vient d'écrire
%% dans un bloc rechercheCARE, l'environnement que se partagent les assistants (Claude,
%% Codex). Dans le PDF, elle dit qui a écrit quoi. Qui répond à un bloc écrit dessous et
%% signe à son tour, sans réécrire ce qui précède. Poussée à droite, ou à la ligne suivante
%% faute de place ; disparaît en mode DIFFUSIONFINALE, comme le bloc qui la porte.
\newcommand\signeCARE[2][]{%
\ifnotmodeDIFFUSIONFINALE{%
{\unskip\nobreak\hfil\penalty50\hskip1em\hbox{}\nobreak\hfil
{\footnotesize\sffamily\bfseries [#2\if\relax\detokenize{#1}\relax\else, #1\fi]}%
\parfillskip=0pt \finalhyphendemerits=0 \par}}}

%% \ARETENIR : la même boîte sous un nom neutre, titre par défaut « À retenir ». C'est le
%% nom employé dans INF412 (décision d'Olivier, 27 septembre 2026), dont les fragments
%% portent un \providecommand de repli pour les tiers, qui n'ont pas ce fichier.
\newcommand\ARETENIR[2][À retenir]{\IMPORTANTCARE[#1]{#2}}
